\chapter{Derivations for the Rotating Ring}
\label{app:ring}

The mechanics behind Chapter~\ref{ch:ring}. The chapter states the gauge
reduction, the exact surrogate score and their consequences, which is what the
argument needs; the algebra establishing them is here, where it can be checked
without interrupting that argument.

\section{Rotations preserve the standard isotropic Gaussian}

Proof of Lemma~\ref{lem:ring-isotropy}.

\begin{proof}
A linear map of a Gaussian is Gaussian, so $R\eta$ is fixed by its first two
moments: $\E[R\eta] = R\,\E[\eta] = 0$ and
$\mathrm{Cov}(R\eta) = R\, I_2\, R^{\T} = R R^{\T} = I_2$. Equivalently, in
terms of densities, $(2\pi)^{-1}e^{-\norm{z}^2/2}$ depends on $z$ only through
$\norm{z}$, and the change of variables $z \mapsto Rz$ changes neither
$\norm{R^{-1}z} = \norm{z}$ nor $|\det R^{-1}| = 1$: the level sets are circles
centred at the origin and a rotation maps each circle to itself.
\end{proof}

\section{The gauge reduction}

Proof of Lemma~\ref{lem:ring-gauge}.

\begin{proof}
Substituting \eqref{eq:ring-sur-dynamics} and using the composition rule of
\eqref{eq:ring-rot-props},
\begin{equation*}
  y_{k+1}
  = R_{-(k+1)\psi}\big(R_\psi z_k + \sigma\eta_{k+1}\big)
  = \underbrace{R_{-(k+1)\psi}R_\psi}_{= R_{-k\psi}} z_k
    + \sigma R_{-(k+1)\psi}\eta_{k+1}
  = y_k + \sigma\widetilde\eta_{k+1} ,
\end{equation*}
and by Lemma~\ref{lem:ring-isotropy} each
$\widetilde\eta_{k+1} \sim \Nd(0,I_2)$; they stay independent because each is a
fixed linear map of a different independent $\eta$. Next, $U_\psi$ is block
diagonal with orthogonal blocks, so
$U_\psi^{\T} U_\psi = \mathrm{diag}(I_2, \dots, I_2) = I_{2L}$ and
$\det U_\psi = \prod_k \det R_{-k\psi} = 1$. Applying $U_\psi$ to
\eqref{eq:ring-channel} gives
$U_\psi x = e^{-t} U_\psi a + \sqrt{\Deltat}\, U_\psi \xi$ with
$U_\psi \xi \sim \Nd(0, I_{2L})$, so $U_\psi x$ has exactly the law of the
noised non-rotating trajectory: the de-rotation commutes with the channel. For
the densities, $z = U_\psi^{\T} y$ and
$P_t(z\,|\,\psi) = P^{(0)}_t(U_\psi z)\,|\det U_\psi|$, the Jacobian factor
being one. Differentiating that identity componentwise, with
$\partial y_b/\partial z_c = (U_\psi)_{bc}$,
\begin{equation*}
  \frac{\partial}{\partial z_c} \log P^{(0)}_t(U_\psi z)
  = \sum_b (U_\psi)_{bc}\,
    \big(\nabla \log P^{(0)}_t\big)_b (U_\psi z)
  = \big[ U_\psi^{\T} \score^{(0)}(U_\psi z, t) \big]_c .
\end{equation*}
The transpose is not cosmetic: a score is a gradient, so it transforms with
$U_\psi^{\T}$ while a point transforms with $U_\psi$.
\end{proof}

\section{The clean law after the gauge}

Iterating \eqref{eq:ring-gauged-rw} and letting the intermediate terms cancel,
\begin{equation}
  y_k - y_0 = \sum_{r=1}^{k} (y_r - y_{r-1})
            = \sigma \sum_{r=1}^{k} \widetilde\eta_r ,
  \qquad\text{so}\qquad
  y_k = A + \sigma \sum_{r=1}^{k} \widetilde\eta_r ,
  \label{eq:ring-unrolled}
\end{equation}
with $A := y_0 \sim p_\lambda$ and the empty sum equal to zero. Every frame
contains the \emph{same} random anchor, drawn once on the ring, plus a partial
sum of independent innovations; the anchor is the only non-Gaussian ingredient
of the model.

The clean law factorises over the ring prior and the one-step transitions, so
$\log P_0(y) = -(\norm{y_0}-1)^2/2\lambda
- (2\sigma^2)^{-1}\sum_{k=0}^{L-2}\norm{y_{k+1}-y_k}^2 + \text{const}$.
Differentiating is a counting exercise: an interior $y_k$ appears in exactly two
summands, the edge arriving at it and the edge leaving it, whose derivatives are
$-(y_k - y_{k-1})/\sigma^2$ and $+(y_{k+1} - y_k)/\sigma^2$. Adding them, and
using $\nabla_y \norm{y} = y/\norm{y}$ for the ring term,
\begin{align}
  \score_0(y,0)
    &= \frac{1-\norm{y_0}}{\lambda}\,\frac{y_0}{\norm{y_0}}
       + \frac{y_1 - y_0}{\sigma^2},
  \notag \\
  \score_k(y,0)
    &= \frac{y_{k-1} - 2y_k + y_{k+1}}{\sigma^2},
       \qquad 1 \le k \le L-2,
  \label{eq:ring-clean-score} \\
  \score_{L-1}(y,0)
    &= \frac{y_{L-2}-y_{L-1}}{\sigma^2} .
  \notag
\end{align} The interior clean score is a
discrete second derivative along internal time, that is a discrete Laplacian:
each block measures how far frame $k$ departs from the average of its
neighbours. It is strictly local, and it vanishes identically on the model's
own noiseless solution, because a trajectory rotating at radius exactly one
maps under the gauge to a constant $y_k$, for which every second difference
vanishes and, with $\norm{y_0} = 1$, so does the ring term. That vanishing is
verified to machine zero (Table~\ref{tab:ring-validation}).

\section{The exact surrogate score}

Proof of Proposition~\ref{prop:ring-exact}.

\begin{proof}
Expanding the exponent of \eqref{eq:ring-Ct} in powers of $\alpha$, with the
correlated covariance $C_t \otimes I_2$, gives
\eqref{eq:ring-anchor-posterior}. Writing
$\log P_t(x) = \text{const} - \tfrac12 x^{\T}(C_t^{-1}\otimes I_2)x
+ \log Z(h(x))$, the gradient of the first term is $-(C_t^{-1}x)_k$ in block
$k$, while $\partial h/\partial x_k = e^{-t}q_k I_2$ and
$\nabla_h \log Z = \E[A\,|\,x]$, the first cumulant identity for the
exponential family in $h$. Differentiating once more, so that the second
derivative of $\log Z$ returns $\mathrm{Cov}(A\,|\,x)$, gives the response.
\end{proof}
